\section{\bt{Design of \system}{\system\ 的设计}}
\label{sec:model}
\bi{\system\ addresses the three challenges through one representation and its lifecycle. Structured definitions separate reusable computation from task parameters. Derived conditions make the data assumptions explicit and connect metric validity to the consuming snapshot. Condition results, repairs, and optimizations are shared across consumers, amortizing the work needed to sustain reuse. The same representation supports entity-level features with per-instance temporal bounds.}
{\system\ 通过统一的表示及其生命周期解决上述三个挑战。结构化定义分离可复用计算与任务参数；派生条件显式表达数据前提，将指标有效性与消费快照联系起来；条件结果、修复与优化在使用者之间共享，分摊持续复用所需的工作。同一表示还支持具有逐样本时间边界的实体级特征。}

\subsection{\bt{What agents share and how it is learned}{智能体共享什么、如何学来}}
\label{sec:shared-object}
\bi{\emph{Shared objects.} \system\ shares database knowledge that agents discover through exploration and whose structural assumptions can be checked against data: metric definitions, including their grain, joins, filters, and encoded value meanings. It also shares table descriptions, verified join paths, and check results to avoid repeated database work. Question parameters and answers remain separate; a result cache handles exact query repeats. Agent strategies, conversations, and preferences remain in agent memory.}
{\emph{共享对象。}\system\ 共享由智能体探索发现、其结构前提可由数据检查的数据库知识：指标定义，包括粒度、连接、过滤及其编码的取值含义。它也共享表描述、验证过的连接路径与检查结果，避免重复的数据库工作。题目参数与答案不属于定义；结果缓存处理完全相同的查询。智能体策略、对话与偏好仍保存在智能体记忆中。}

\bi{A definition revision is $m^r=(B,I,C,E,r)$, where $B$ is its business meaning, $I$ its structured implementation, $C$ its set of executable conditions, $E$ its learning record, and $r$ its revision number. The implementation records the fact table, aggregate expression, grain, date join, joins, and definition filters; its joins form a many-to-one tree rooted at the fact table (a star or snowflake). Separating question parameters from $I$ enables compilation for new periods, differences, and rankings. Keeping $C$ explicit makes validity a relation between the definition and the database snapshot: the same revision can remain usable through many updates that preserve its conditions. Figure~\ref{fig:metric-record} shows how this separation supports both reuse and maintenance.}
{定义的一个版本记为 $m^r=(B,I,C,E,r)$，其中 $B$ 为业务含义，$I$ 为结构化实现，$C$ 为可执行条件集，$E$ 为学习记录，$r$ 为版本号。实现记录事实表、聚合表达式、粒度、日期连接、连接和定义过滤；连接构成以事实表为根的多对一连接树（星型或雪花型）。将问题参数与 $I$ 分离，使定义能够编译为新期间、差值和排名所需的查询。显式保留 $C$，使有效性成为定义与数据库快照之间的关系：同一版本可在多次保持其条件成立的更新后继续使用。图~\ref{fig:metric-record} 展示这一分离如何同时支持复用与维护。}

\bi{\emph{Extraction and publication.} A candidate comes from a session with a stated business definition and an independently verified answer, supplied by a user, report, or analyst; in the benchmark evaluation, an automated oracle compares the final scalar answer with the reference answer on the same snapshot before extraction. The LLM extractor receives the question, SQL and arithmetic trace, rejected queries, and schema and join metadata, but no business rows or result values. It proposes a JSON record; the server parses its fields and derives evidence and conditions itself. Extraction is therefore model-based, whereas publication uses deterministic checks: (G1) an explicit business-definition basis; (G2) an independently judged answer and verified calculation chain; (G3) valid tables, columns, joins, expressions, and definition filters, with question parameters excluded; (G4) fact-grain uniqueness; (G5) a single read-only example query that passes execution review; (G6) unchanged learning dependencies and reproduction of the verified value by the example; and (G7) reproduction of that value by SQL compiled from the structured fields. Numeric comparisons use the learning task's declared precision. G7 tests the artifact consumers will execute, catching omissions introduced during extraction.}
{\emph{提取与发布。}候选来自题目明确业务定义、答案经过独立核实的会话；核实可来自用户、报表或分析师，基准评估中，自动判题器在提取前将最终标量答案与同一快照上的参考答案比较。LLM 提取器接收题目、SQL 与算术轨迹、被拒绝的查询以及模式和连接元数据，但看不到业务行或结果值。它提出 JSON 记录；服务器解析字段，自行推导证据与条件。因此提取使用模型，发布则执行确定性检查：（G1）有明确的业务定义依据；（G2）答案经过独立判定且计算链可验证；（G3）表、列、连接、表达式与定义过滤合法，排除题目参数；（G4）事实粒度唯一；（G5）示例是通过执行审查的单条只读查询；（G6）学习依赖未变化，示例重现核实值；（G7）由结构化字段编译的 SQL 重现该值。数值比较使用学习任务声明的精度。G7 检查消费者实际执行的工件，发现提取时遗漏的条件。}

\bi{\emph{Use.} An agent queries \code{find\_metric} with a metric name or phrase, matched against canonical names and aliases extracted during learning. The agent receives the valid revision: its business definition, structured fields, generated SQL, and revision number. It sets its own period in the generated SQL, or wraps the SQL for a difference or a ranking, declares the revision to \code{run\_sql}, and answers. This replaces the exploration: the agent needs no table listing, value probe, or join test. On TPC-DS, such a session takes three LLM calls for a metric whose definition needs three inner joins; agents without memory spend \PpNoneMTwoSteps\ calls on questions about it, because they must find the fact table, the measure, and every join themselves. Agents issue at most three memory lookups per question before falling back to standard schema exploration.}
{\emph{使用。}智能体向 \code{find\_metric} 提交指标名或短语，与学习过程中提取的规范名称及别名匹配。智能体得到有效版本：业务定义、结构化字段、生成的 SQL 与版本号。它在生成的 SQL 中填入自己的期间，或在外面包一层求差值或排名，在 \code{run\_sql} 中声明该版本，然后作答。这就取代了探索：智能体不必列表、探查取值或检验连接。在 TPC-DS 上，对一个定义需要三个内连接的指标，这样的会话只用三次大模型调用；无记忆的智能体在关于它的题上要花 \PpNoneMTwoSteps\ 次调用，因为事实表、度量和每个连接都要自己找。智能体每道题最多查找三次记忆，未命中时转入常规模式探索。}

\subsection{\bt{When shared knowledge is correct}{共享的知识何时正确}}
\label{sec:sufficiency}
\bi{Data-only updates break learned SQL while the schema stays the same: status-change rows, restatements, or duplicate loads violate the grain, reused dimension keys make a many-to-one join fan out, and late or reformatted date keys move facts to another period or to none~\cite{kimball2013toolkit}. A condition $c$ is a predicate over the tables $\operatorname{dep}(c)$; its check query $Q_c$ returns statistics used to evaluate the predicate. Its normalized identity $\operatorname{id}(c)$ includes the condition type, tables, keys, join direction, filters, and null policy. \system\ derives four kinds from $I$. \emph{Filtered key uniqueness} $\operatorname{KeyUnique}(T,K,\varphi)$ checks uniqueness among rows of $\sigma_\varphi(T)$ with non-null key components. \emph{Join cardinality} limits matches for each left row. \emph{Date join} checks the fact's date key against the date dimension. \emph{Join completeness} bounds the unmatched fraction for the date join and each inner join by its publication baseline plus 0.1 percentage points; path checks count facts reaching the final edge.}
{纯数据更新会在模式不变的情况下破坏学到的 SQL：状态流水、更正保留旧行或重复加载破坏粒度，被重用的维度键使多对一连接放大行数，迟到或改了格式的日期键把事实移到别的期间或使其无处可归~\cite{kimball2013toolkit}。条件 $c$ 是表集合 $\operatorname{dep}(c)$ 上的谓词，其检查查询 $Q_c$ 返回用于判定条件的统计量；规范化标识 $\operatorname{id}(c)$ 包含条件类型、表、键、连接方向、过滤和空值规则。\system\ 从 $I$ 推导四类条件。\emph{过滤后键唯一性} $\operatorname{KeyUnique}(T,K,\varphi)$ 检查 $\sigma_\varphi(T)$ 中各键分量非空的行是否唯一。\emph{连接基数}限制每条左侧行的匹配数。\emph{日期连接}检查事实日期键与日期维度的连接。\emph{连接完整性}要求日期连接和每个内连接的未匹配比例不超过发布时基线加 0.1 个百分点；路径检查按到达末端连接的事实行计数。}

\bi{For a period $p$, the generated SQL evaluates $\mathrm{AGG}(e)$ over the rows of $\sigma_\varphi(F)\bowtie R_1\bowtie\cdots\bowtie R_k$ whose date falls in $p$. The business definition supplies three domain premises: (B1) the selected fact rows have non-null keys $K$ whose values correspond one-to-one to business events; (B2) $e$ evaluates to each event's measurement; and (B3) the date assignment and definition filters represent the intended business selection.}
{对期间 $p$，生成的 SQL 在 $\sigma_\varphi(F)\bowtie R_1\bowtie\cdots\bowtie R_k$ 中日期落在 $p$ 内的行上计算 $\mathrm{AGG}(e)$。业务定义提供三条领域前提：（B1）选中的事实行具有非空键 $K$，其取值与业务事件一一对应；（B2）$e$ 求值得到每个事件的度量值；（B3）日期归属与定义过滤表达预期的业务选择。}
\begin{proposition}[\bt{Sufficiency}{充分性}]
\label{prop:sufficiency}
\bi{If B1--B3 and $C(m^r)$ hold on $D_s$, then for any period $p$ each retained qualifying event appears exactly once in the input to $\mathrm{AGG}(e)$, and the unmatched fraction on each checked table or path is bounded by its publication baseline plus 0.1 percentage points. When all required date and inner joins match completely, the aggregate is exact for $p$.}
{若 B1--B3 与 $C(m^r)$ 在 $D_s$ 上成立，则对任意期间 $p$，每个保留下来的合格事件在 $\mathrm{AGG}(e)$ 的输入中恰好出现一次，且每个受检查表或路径的未匹配比例不超过发布时基线加 0.1 个百分点。若全部所需日期连接和内连接都完整匹配，则该期间的聚合结果精确。}
\end{proposition}
\begin{proof}[\bt{Proof sketch}{证明概要}]
\bi{B1 and filtered key uniqueness give one fact row per event. Join cardinality prevents duplication; date and inner joins may remove unmatched rows. B2 and B3 assign each retained row its measurement and intended period. Complete matching retains all qualifying events, while join-completeness predicates bound the unmatched fraction on their respective tables and paths.\ifbilingual\else\qedhere\fi}
{B1 与过滤后键唯一性保证每个事件有一条事实行。连接基数防止行被复制；日期连接与内连接可能移除未匹配的行。B2、B3 将每条保留的行归入预期期间并提供其度量值。完整匹配保留全部合格事件；连接完整性谓词则约束各受检查表和路径的未匹配比例。\qedhere}
\end{proof}

\bi{Table- and path-level completeness checks establish population-level bounds for the checked tables and paths, allowing period queries to share the same validation results.}
{表级和路径级完整性检查为受检查的表与路径建立总体层面的界限，使不同期间查询能够共享同一组验证结果。}

\subsection{\bt{Maintenance once for all agents}{为所有智能体一次性维护}}
\label{sec:maintenance}
\bi{\emph{Shared check results.} The version $V(T)$ of a table identifies the state that a check read: either the schema fingerprint, DML counters, and ETL batch identifier, or a \emph{transactional version} that a statement-level trigger increments inside every writing transaction. A check result is keyed by $\bigl(\operatorname{id}(c),V(\operatorname{dep}(c))\bigr)$, so it serves every definition and every agent that needs the same condition on the same table versions. When an agent looks up a definition, \system\ compares the versions recorded for each condition's tables with the current ones, keeps the check results of conditions over unchanged tables, looks up a completed result for each remaining condition at the current versions, and runs $Q_c$ only if none exists; concurrent requests for the same check share that one execution (request coalescing). Joins are checked through shared join entries, so a changed or vanished join path also makes the definition pending. A failed condition invalidates the revision; consumers then learn that the definition is unavailable and explore on their own, and the next successful learning session publishes it again.}
{\emph{共享的检查结果。}表的版本 $V(T)$ 标识检查读到的状态：或者组合结构指纹、DML 计数与 ETL 批次号，或者使用\emph{事务性版本}，由语句级触发器在每个写事务内递增。检查结果以 $\bigl(\operatorname{id}(c),V(\operatorname{dep}(c))\bigr)$ 为键，因此服务于所有在相同表版本上需要同一条件的定义和智能体。智能体查找定义时，\system\ 把每个条件所读各表记录下的版本与当前版本比较：未变化表上的条件保留原有检查结果；对其余条件，查找同一规范化形式在当前版本上的已完成结果，没有才执行 $Q_c$，同一检查的并发请求共享这一次执行（请求合并）。连接通过共享的连接条目检查，因此连接路径被改变或消失时，定义同样变为待验证。条件失败使版本失效；使用者随即得知定义不可用、自行探索，下一次成功的学习会话会重新发布它。}

\bi{\emph{Update cost and scheduling.} In the transactional-version implementation, a statement-level trigger performs one metadata upsert per modified table and writing statement, distributed over 16 counter shards to reduce write contention. Validation runs lazily on demand: after an update, the first consumer triggers and completes any uncached dependency checks before use; concurrent consumers coalesce onto that execution, and subsequent consumers reuse its results while dependency versions match.}
{\emph{更新代价与调度。}事务性版本实现中，语句级触发器对每个被修改的表、每条写语句执行一次元数据 upsert，分散到 16 个计数分片以减少写争用。验证采用按需惰性执行：更新后，首个消费者在使用前触发并完成未缓存的依赖检查；并发消费者合并到同次执行，后续消费者在依赖版本一致时复用其结果。}

\bi{\emph{Repair once.} When a fact-table grain fails, deterministic repair searches equality filters on columns with few distinct values that restore one row per business key without losing keys. Publication requires a unique passing filter to prevent ambiguous matches, such as competing primary and backup copies, and reproduction of the learning query's result through an \emph{as-of query} on the learning snapshot. A repair runs once and serves every agent; requests that arrive while it runs join it, and a definition that cannot be repaired is invalidated, so its consumers explore until a session relearns it.}
{\emph{只修复一次。}事实表粒度失效时，修复在取值较少的列上搜索等值过滤，要求恢复每个业务键一行且不丢键。发布要求恰好一个过滤通过，以避免主数据与备份副本等歧义匹配；同时，候选过滤须通过学习快照上的 \emph{as-of 查询} 重现学习查询的结果。修复只做一次、服务所有智能体；修复期间到达的请求与之合并；无法修复的定义被置为失效，其使用者自行探索，直到某个会话重新学到它。}

\bi{\emph{Checks at query execution.}\label{sec:execution} An agent declares the revisions a query uses, and \code{run\_sql} checks these revisions and required grain filters. For a declared revision whose metadata remains stable during condition selection, snapshot-bound execution opens a repeatable-read transaction, retrieves condition results at the snapshot's transactional versions, and runs missing checks before executing the query on that snapshot. Proposition~\ref{prop:sufficiency} establishes correctness conditions for compiled aggregates used by period and ranking queries. Manually altered SQL undergoes read-only safety review. The following lemma establishes when condition results can be reused across snapshots.}
{\emph{执行时检查。}智能体声明查询所用的版本，\code{run\_sql} 检查这些版本及粒度所需过滤。对于在条件选择期间元数据保持稳定的声明版本，快照绑定执行开启可重复读事务，按快照的事务性版本取得条件结果，运行缺失的检查，再在同一快照上执行查询。命题~\ref{prop:sufficiency} 确立期间与排名查询所用编译聚合的正确性条件。手动修改的 SQL 接受只读安全审查。下面的引理给出跨快照复用条件结果的依据。}

\label{sec:snapshot}
\begin{lemma}[\bt{Check-result reuse}{检查结果复用}]
\label{lem:reuse}
\bi{For read-only snapshots under the write-coverage assumptions below, let $v_s(T)$ be the transactional version of $T$ visible in snapshot $s$. If $v_{s_1}(T)=v_{s_2}(T)$ for every $T\in\operatorname{dep}(c)$, then $c(D_{s_1})=c(D_{s_2})$.}
{对于满足下述写入覆盖前提的只读快照，设 $v_s(T)$ 为快照 $s$ 中可见的 $T$ 的事务性版本。若对每个 $T\in\operatorname{dep}(c)$ 都有 $v_{s_1}(T)=v_{s_2}(T)$，则 $c(D_{s_1})=c(D_{s_2})$。}
\end{lemma}
\begin{proof}
\bi{The visible committed-transaction sets are nested in snapshot acquisition order, say $s_1\subseteq s_2$. A transaction that writes $T$ increments $T$'s counter in the same transaction, so equal versions leave no transaction in $s_2\setminus s_1$ that writes $T$, and both snapshots see the same state of every table that $Q_c$ reads.}
{按快照获取顺序，可见的已提交事务集合相互包含，不妨设 $s_1\subseteq s_2$。写 $T$ 的事务在同一事务内递增 $T$ 的计数，因此版本相等意味着 $s_2\setminus s_1$ 中没有写 $T$ 的事务，两个快照看到 $Q_c$ 所读各表的同一状态。}
\end{proof}
% PostgreSQL 18: sql-copy.html (triggers/FREEZE), sql-createtrigger.html
% (foreign-key actions), sql-truncate.html (transactional but not MVCC-safe).
\bi{\emph{Operational assumptions.} Statement-level triggers track inserts, updates, and deletes (including \code{COPY FROM} and cascading foreign-key actions) on unpartitioned base tables without inheritance. Version increments commit or roll back atomically with writes, and counters are never reset or reused while cached checks remain valid. Non-MVCC operations, such as \code{TRUNCATE}, table rewrites, and \code{COPY FREEZE}, require quiescent readers. The deployment invalidates affected check results before reuse after DDL, trigger-bypassing writes, or restore/reset operations.}
{\emph{运行假设。}语句级触发器跟踪无分区、无继承的基表上的插入、更新和删除（包括 \code{COPY FROM} 和外键级联操作）。版本递增与写入原子地提交或回滚；缓存检查有效期间不重置或复用计数器。非 MVCC 操作（如 \code{TRUNCATE}、表重写及 \code{COPY FREEZE}）要求没有在途读事务。发生 DDL、绕过触发器的写入或恢复／重置操作后，部署须在复用前使受影响的检查结果失效。}

\subsection{\bt{Optimization once for all agents}{为所有智能体一次性优化}}
\label{sec:optimization}
\bi{Optimizing a published definition reduces execution cost for every consumer of its generated SQL. After learning, \system\ runs an optimization pass over the valid definitions; candidates are structured definitions, so compilation, the conditions, and Proposition~\ref{prop:sufficiency} apply to them. Two rules propose candidates: \emph{key range} replaces the period join with a range predicate on the fact table's date key, whose bounds scalar subqueries read from the date dimension, and \emph{left-join removal} drops a left join that contributes no column and no filter. The optimizer first evaluates these provably equivalent rules (Proposition~\ref{prop:rewrite}). If no rule candidate is adopted, an optional LLM advisor analyzes the definition, generated SQL, execution plan, and table sizes to propose alternatives. \system\ publishes validated candidates with a rule-supported structure and retains other passing proposals as suggestions for review. A candidate must pass the static, grain, and SQL checks and three tests. \emph{Assumption}: the key-range rule adds date-key contiguity, checked before adoption and maintained afterwards. \emph{Equivalence}: candidate and current revision return equal results for the learning question as of learning time and, within one repeatable-read snapshot, for the twelve months of the learning year, a quarter, a period difference, and a monthly ranking. \emph{Cost}: in the same snapshot, the paired execution times of at least eight periods, measured in alternating order, must have a 95\% confidence interval below zero and a mean saving of at least 10\%.}
{优化一个已发布的定义，可以降低所有使用者执行其生成 SQL 的代价。学习结束后，\system\ 对有效的定义执行一轮优化；候选都是结构化定义，因此编译、条件和命题~\ref{prop:sufficiency} 都适用于它们。两条规则提出候选：\emph{日期键范围}把期间连接改为事实表日期键上的范围谓词，上下界由标量子查询从日期维度读出；\emph{去掉左连接}去掉一个既不提供列也不带过滤的左连接。优化器首先评估这两条可证明等价的规则（命题~\ref{prop:rewrite}）。若未采纳规则候选，可选的 LLM 顾问分析定义、生成的 SQL、执行计划与表规模并提出备选方案。\system\ 发布通过验证且结构受规则支持的候选，其他通过检查的提议则保存为待审阅建议。候选必须通过静态、粒度和 SQL 检查以及三项测试。\emph{前提}：日期键范围规则新增日期键连续条件，采纳前检查，之后受维护。\emph{等价}：候选与当前版本对学习题按学习时刻返回相同结果，并在一个可重复读快照内对学习年份的 12 个月、一个季度、一次期间差值和一次按月排名返回相同结果。\emph{代价}：在同一快照中，至少 8 个期间以交替顺序测得的配对执行时间差，其 95\% 置信区间必须整体低于 0，平均节省至少 10\%。}
\begin{proposition}[\bt{Rule equivalence}{规则等价}]
\label{prop:rewrite}
\bi{Let the date dimension's key be unique on $D$. (a) If \emph{date-key contiguity} holds on $D$, i.e., the keys of every month form an integer interval that begins one above the previous month's last key, then for every period of consecutive months the generated SQL of a key-range revision returns on $D$ what that of its previous revision returns. (b) If the right side of a removed left join is unique on the join key on $D$, removing the join does not change the result.}
{设日期维度的键在 $D$ 上唯一。（a）若\emph{日期键连续性}在 $D$ 上成立，即每个月的键构成一个整数区间、且从上个月最后一个键加一开始，则对由连续月份组成的每个期间，键区间版本的生成的 SQL 在 $D$ 上返回与其前一版本相同的结果。（b）若被消除的左连接右侧在连接键上唯一，消除该连接不改变结果。}
\end{proposition}
\begin{proof}[\bt{Proof sketch}{证明概要}]
\bi{(a) By contiguity, the keys of a period's months are exactly the integer interval between the bounds, so a fact row's key lies in the interval if and only if it joins one date row of the period. (b) A left join with at most one match keeps each left row once and supplies nothing that is read.}
{（a）由连续性，一个期间各月的键恰好是边界之间的整数区间，因此事实行的键落在区间内当且仅当它连接到该期间的一行日期。（b）至多一个匹配的左连接使每个左侧行恰好保留一次，且不提供任何被读取的内容。}
\end{proof}
\bi{The assumption of (a) is a condition of the new revision; that of (b) is a condition of the previous revision that the new revision no longer needs. Wherever the maintained conditions hold, a rule revision thus agrees with its previous revision on every period, and the sampled tests check the implementation and measure the cost. Maintenance checks an added assumption last; when it alone fails, every condition shared with the previous revision has just been confirmed, and \system\ invalidates the optimized revision and makes its previous revision current again.}
{（a）的前提是新版本的条件；（b）的前提是前一版本的条件，新版本不再需要它。因此，只要受维护的条件成立，规则版本就在每个期间上都与前一版本一致；抽样测试检查的是实现并测量代价。维护最后检查新增的前提；只有它不成立时，与前一版本共有的每个条件都刚刚得到确认，\system\ 使优化后的版本失效，并让前一版本重新成为当前版本。}

\subsection{\bt{Prediction tasks: shared feature definitions}{预测任务：共享的特征定义}}
\label{sec:features}
\bi{Data agents also build entity-level feature tables for predictive modeling. Given seed rows with an entity and a seed time, feature SQL returns aggregates of the entity's data before that time, and a fixed learner trains on these features. RelBench's forecasting tasks follow this form~\cite{robinson2024relbench}. A \emph{feature definition} extends a metric definition (fact, measure, grain, verified many-to-one joins, and required filters) with three fields: the entity reached through a foreign key or verified join, optionally through an anchor column such as a seed post's owner; a look-back window; and the fill value for a seed with no matching rows. Compilation produces one lateral aggregate per seed and bounds the definition's event-time column by the seed time and window.}
{数据智能体也为预测建模构建实体级特征表。给定包含实体和种子时间的种子行，特征 SQL 返回该实体在种子时间之前的数据聚合，再由固定学习器训练。RelBench 的预测任务采用这一形式~\cite{robinson2024relbench}。一条\emph{特征定义}在指标定义（事实表、度量、粒度、经验证的多对一连接及必需过滤）上增加三个字段：经外键或已验证连接到达的实体，必要时经锚列（如种子帖子的作者）定位；回看窗口；没有匹配行时的填充值。编译为每个种子生成一个 LATERAL 聚合，并根据种子时间与窗口约束定义中的事件时间列。}

\bi{\emph{Two-tier reuse.} \system\ pairs feature definitions with \emph{task programs}: the definitions submitted by a successful session, its validation score, and its complete SQL. A task program provides an executable starting point for retraining the same task; individual definitions support selection, composition, and refinement across compatible tasks sharing an entity. Complete task SQL preserves expressions such as self-joins, while structured definitions expose reusable feature components. Consumer sessions validate candidate SQL on the current task and audit temporal leakage before adoption.}
{\emph{两层复用。}\system\ 将特征定义与\emph{任务方案}结合；任务方案记录成功会话提交的定义、验证分数和完整 SQL。任务方案为同一任务的重训提供可执行起点；单条定义支持在共享实体的兼容任务间选择、组合和细化特征。完整任务 SQL 保留自连接等表达式，结构化定义则提供可复用的特征组件。消费者会话在当前任务上验证候选 SQL，并在采用前审计时间泄漏。}

\bi{\emph{Publication.} After a session whose features beat the entity-table reference, an extractor turns the submitted SQL into definition drafts. Publication checks four properties on the current snapshot: valid fields (existing columns, expressions confined to the involved tables, and many-to-one joins); uniqueness of the fact's grain key; value-by-value reproduction of the submitted column on 2,000 seeds by the compiled SQL; and unchanged values when recomputed on views that hide rows after the seed time. Drafts that pass these checks enter the shared library.}
{\emph{发布。}一次会话的特征优于实体表参照时，提炼器将提交的 SQL 转为定义草稿。发布前在当前快照上检查四项性质：字段合法（列存在、表达式只读涉及的表、连接为多对一）；事实表的粒度键唯一；编译 SQL 在 2,000 个种子上逐值重现提交的列；在隐去种子时间之后的行的视图上重算，值保持不变。通过检查的草稿进入共享库。}

\bi{\emph{Retrieval and maintenance.} \code{find\_program} returns the task's latest program: complete SQL checked for executability on the current snapshot and SQL compiled from its valid definitions. \code{find\_features} retrieves definitions published by any agent for the task's entity and compiles them for the consumer's seed table; \code{compose\_features} compiles a selected set. After data updates, the shared maintenance protocol rechecks affected conditions and revokes definitions whose conditions fail. Subsequent successful sessions replenish the library through the same publication checks.}
{\emph{检索与维护。}\code{find\_program} 返回任务最近的方案，包括已检查在当前快照上可执行的完整 SQL，以及由有效定义编译的 SQL。\code{find\_features} 检索任意智能体为任务实体发布的定义，并针对消费者的种子表编译；\code{compose\_features} 编译选定集合。数据更新后，共享维护协议重验受影响的条件，撤销条件失败的定义。后续成功会话通过同一套发布检查补充共享库。}
